% II. Анализ на възможните решения и обосновка на избора.
\section{Анализ на възможните решения и обосновка на избора}
\label{sec:alternatives}

Всяка от трите подзадачи допуска повече от едно решение и изборът между тях предопределя
както обхвата, така и това, което изобщо може да бъде измерено. Разделът излага
разгледаните възможности и мотивите за приетото решение. Където реализацията впоследствие
наложи различен избор от първоначално предвидения, това е записано тук заедно с причината.

\subsection{Изисквания към сглобяването и тяхната тежест}
\label{subsec:dependencies}

Преди трите подзадачи стои едно решение, което ги предопределя и трите: от какво има право
да зависи сборката. Разгледани са два подхода.

Първият е обичайният: необходимите библиотеки се обявяват и се доставят от управител на
пакети или се изтеглят при настройването на сборката. Реализацията е по-кратка, защото
готовите библиотеки поемат работа, но сборката изисква подготвена машина и мрежа.

Вторият е сборката да зависи единствено от компилатор за C++20 и CMake. Тогава всяка
функция, която иначе би дошла отвън, трябва или да бъде написана, или да отпадне заедно с
това, което зависи от нея.

Приема се вторият подход. Основанието е, че хранилището е публично и трябва да може да
бъде сглобено от непознат без предварителна подготовка. Сборка, която изисква управител на
пакети, е сборка, която никой не пуска. Приетото решение има цена и тя е платена на три
места, изброени в Раздел~\ref{subsec:build}: отпада библиотеката за прихващане на пакети,
отпада външната рамка за тестове и отпада проверката на подписа. Първите две отпадания не
отнемат нищо от задачата. Третото отнема и затова е разгледано отделно в
Раздел~\ref{subsec:scope}.

\subsection{Източник на данните за анализа}
\label{subsec:source}

Първата възможност е анализът да се извършва вътре в приложението, чрез обратни извиквания
на криптографската библиотека, която вече е разчела съобщенията. Реализацията тогава е
кратка, но резултатът описва само собствените връзки на приложението и е обвързан с
версията на библиотеката.

Втората възможност е пасивно наблюдение на трафика чрез библиотека за прихващане на пакети,
върху жив интерфейс или върху записан файл. Анализът обхваща чужди на приложението връзки,
а записаният файл прави измерването възпроизводимо. Цената е една задължителна външна
зависимост, което влиза в противоречие с Раздел~\ref{subsec:dependencies}.

Третата възможност, която е и приетата, е слоят за разчитане да приема просто изглед към
байтове, тоест \code{std::span<const std::uint8\_t>}. Откъде идват байтовете не е негова
работа: записан файл, сокет или генериран от тестовете набор изглеждат еднакво. Тази
възможност запазва двете предимства на втората, защото анализът пак не е обвързан с
приложението и входът пак е фиксиран, а премахва зависимостта. Загубеното е удобството да
се чете формат на записан трафик направо; това е отделна задача за разчитане на формат, а
не свойство на анализа, и е записано като посока за по-нататъшна работа.

\subsection{Обхват на собствената реализация}
\label{subsec:scope}

Криптографските примитиви са зона, в която собствената реализация е грешка. Изборът на
алгоритми, устойчивостта им спрямо атаки по странични канали и правилността на самите
изчисления са резултат от години рецензии. Първоначалният замисъл затова предвиждаше
примитивите да се вземат от утвърдена библиотека.

При реализацията това изискване влезе в противоречие с
Раздел~\ref{subsec:dependencies}, който е по-силен. Разгледани са три изхода.

Първият е да се приеме задължителна криптографска зависимост, което отменя изискването за
сглобяване без подготовка. Вторият е да се напише собствена аритметика с големи числа, което
е точно грешката, срещу която предупреждава първият абзац. Третият е проверката на подписа
да не се извършва и това да бъде обявено при всяка верига, редом с проверките, които са
проведени.

Приема се третият изход. Той е по-малкото зло по следната причина: първите два или отнемат
свойството, заради което хранилището е използваемо, или заменят рецензирана реализация с
нерецензирана. Третият не отнема нищо и не заменя нищо, а само стеснява обхвата и го
обявява. Съществено е, че обявяването е част от изхода, а не добавка към него. Мълчаливото
пропускане на проверка прави отчета по-опасен от липсата на отчет, защото читателят би
приел, че веригата е проверена криптографски.

Логиката над примитивите остава предмет на проекта, както беше предвидено. Разчитането на
съобщенията на ръкостискането, изграждането на пътя до доверен корен, проверката на
ограниченията върху имената и на разширенията, както и обработката на състоянието на
отмяната, се реализират самостоятелно по описанието в \cite{rfc5280}.

Проверката на отмяната е сведена до локално поддържан списък от серийни номера. Изтегляне
на списък с отменени сертификати или запитване по \cite{rfc6960} изискват изходяща мрежова
връзка, каквато системата няма и по устройство не може да има. Политиката при липсваща
информация за отмяна остава явен параметър, както е разгледано в
Раздел~\ref{subsubsec:revocation}.

\subsection{Стратегия за допускане на връзки}
\label{subsec:admission}

Разгледани са четири механизма за защита при лавинен трафик, всеки от които има различна
цена. Криптографски изчислените стойности без съхранение на състояние, описани в
\cite{rfc4987}, премахват нуждата от състояние за незавършена връзка, но изискват
допълнителен обмен. Ограничаването на честотата по източник е евтино, но е безполезно срещу
разпределен източник и вредно, когато честният и лавинният трафик делят един адрес.
Допускането срещу изчислителна задача прехвърля цената върху клиента, но внася закъснение
при всяко първо свързване. Отчитането на таблицата с връзките е по-скоро измерителен,
отколкото защитен механизъм, но без него останалите три не могат да бъдат оценени.

Приема се всичките четири да бъдат реализирани и включвани поотделно, вместо да се избере
един. Причината е, че въпросът на проекта не е кой механизъм е най-добър, а колко струва
всеки от тях. Отговорът изисква сравнение при еднакви условия, тоест изисква всички да
съществуват в една и съща система.

Към четирите се добавя политика на плавно влошаване, която не е самостоятелен механизъм, а
правило, което вдига цената на останалите три с нарастването на запълването на таблицата.
Тя е добавена, защото без нея механизмите работят с една и съща настройка при празен и при
пълен сървър, а точно при пълен сървър защитата има значение.

\subsection{Език и инструментариум}
\label{subsec:language}

Изборът на C++20 е продиктуван от три съображения. Разчитането на протокола работи с
байтови буфери, за което \code{std::span} дава невладеещ изглед без копиране и без загуба
на дължината. Измерването на закъснението изисква предвидимо време на изпълнение, без паузи
от събиране на боклука. И накрая, стандартната библиотека покрива нужното, така че
изискването от Раздел~\ref{subsec:dependencies} е изпълнимо без да се пише всичко наново.

Алтернативите бяха разгледани и отхвърлени: Rust заради по-малкия личен опит на автора в
рамките на срока на проекта, а езиците със среда за изпълнение и събиране на боклука заради
влиянието им върху точно тази величина, която проектът измерва.

За сглобяването е избран CMake, защото е стандартът за този инструментариум и защото
позволява отделните цели за компилация да направят зависимостите между слоевете видими в
конфигурацията, а не само в текста. Тестовете се регистрират в CTest поотделно, като
списъкът им се произвежда от самия изпълним файл след свързването, за да не може да се
размине с кода.
